% Source: EGAIV.1.pdf, PDF pp.68--77; printed pp.69--78.
\setcounter{section}{18}
\section{形式光滑代数与Cohen环}
\source{EGA IV第一部分，本节原文为印刷第69--115页；本轮连续译入第69--89页，至（19.4.7）}
\setcounter{subsection}{-1}
\subsection{引言}
\egap{19.0.1}
在第IV章中，我们将引入并研究一类重要的预概形态射，即\emph{光滑态射}\orig{某些近期著作也称之为“简单态射”，这一用语取自经典的“简单点”术语；但是它容易引起混淆，尤其是在代数群理论中。}。它们的一项基本性质——这一性质与一个有限性条件合起来便刻画了它们——是态射的\emph{提升性质}：若 $f:X\to Y$ 是光滑态射，$g:Y'\to Y$ 是态射，则对于使 $Y''$ 成为与 $Y'$“相差甚微”的预概形的任意态射 $h:Y''\to Y'$，每个 $Y$-态射 $Y''\to X$ 都分解为 $Y''\xrightarrow{h}Y'\to X$。更精确地，先限于 $Y=\Spec(A)$、$X=\Spec(B)$ 均仿射的情形：若对每个 $A$-代数 $C$ 及 $C$ 的每个幂零理想 $\mathfrak j$，任意 $A$-同态 $B\to C/\mathfrak j$ 都可以提升为 $B\to C\to C/\mathfrak j$，则称 $B$ 为\emph{形式光滑的 $A$-代数}。换言之，映射
\[
 \Hom_A(B,C)\longrightarrow\Hom_A(B,C/\mathfrak j)
\]
是满射。在许多应用中，$X$ 将作为某个反变函子 $Y'\mapsto F(Y')$ 的代表对象出现；该函子从 $Y$-预概形范畴取值于集合范畴，因此有 $F(Y')\simeq\Hom_Y(Y',X)$（$0_{\mathrm{III}}$，8.1.8）。在仿射情形，令 $F^\circ(C)=F(\Spec(C))$，则验证上述条件下 $F^\circ(C)\to F^\circ(C/\mathfrak j)$ 为满射——即使尚不知道 $F$ 可表，也能进行这一验证——便可在附加的有限性条件满足时证明 $f$ 光滑。\cmt{F6-H01}{本段开头是引导性的描述，不应脱离后面的仿射条件，读作任意非仿射测试预概形都存在全局提升。正文在（19.3.1）给出拓扑代数的精确定义；几何情形须区分局部提升与全局提升。}

\egap{19.0.2}
对于拓扑环 $A$ 上的拓扑代数，也有相应的形式光滑代数概念；这里暂不作精确定义（参见19.3.1）。§19首先从初等观点研究这些概念；随后借助将在§§20、21中发展的微分性质，在§22证明更深入的定理。这里概述关于形式光滑代数的主要结果。

\noindent I．设 $A,B$ 为两个诺特局部环，$\varphi:A\to B$ 为局部同态，$k$ 为 $A$ 的剩余域，并令 $B_0=B\otimes_Ak$。给 $A,B,B_0$ 配备各自的预进制拓扑，给 $k$ 配备离散拓扑。$B$ 为形式光滑 $A$-代数的充要条件是：$B$ 为平坦 $A$-模，且 $B_0$ 为形式光滑 $k$-代数（19.7.1）。因此，对于诺特局部环，该定理把形式光滑性归结为域上的诺特局部代数的同一问题。

\noindent II．设 $k$ 为域，$A$ 为一个诺特局部 $k$-代数。$A$ 形式光滑的充要条件是它在 $k$ 上\emph{几何正则}，即对于 $k$ 的每个有限扩张 $k'$，半局部环 $A\otimes_kk'$ 都正则（22.5.8）；这时 $A$ 的完备化 $\widehat A$ 同构于形式幂级数环 $K[[T_1,\ldots,T_n]]$（19.6.5）。此外，当 $A$ 完备且在 $k$ 上形式光滑时，它的 $k$-代数结构完全由其剩余域 $K$ 及维数确定；只要满足不等式
\[
 \dim(A)\geq\operatorname{rg}_{K}(\Upsilon_{K/k}),
\]
该维数可以任意，其中 $\Upsilon_{K/k}$ 是 $K$ 的“不完美模”（22.2.6）。\cmt{F6-T01}{初版本处秩的下标印作 $k$；本轮已直接核对同卷（22.2.6）（印刷第188页），该定理明确以 $K$-向量空间的秩表述同一维数界。因此这里明示统一为 $K$。核对该引文不计作§22已经译入。}特别地，$k$ 的一个扩张 $K$ 形式光滑的充要条件是 $K/k$ 可分（19.6.1）。

\noindent III．设 $A$ 为诺特局部环，$\mathfrak j\ne A$ 为其理想，$A_0=A/\mathfrak j$；设 $B_0$ 为完备诺特局部环，$A_0\to B_0$ 为使 $B_0$ 成为形式光滑 $A_0$-代数的局部同态。那么存在完备诺特局部环 $B$、使 $B$ 成为平坦 $A$-模的局部同态 $A\to B$，以及 $A_0$-同构 $B\otimes_AA_0\simeq B_0$（于是由I，$B$ 是形式光滑的 $A$-代数）；此外，这些性质将 $B$ 确定到同构（19.7.2）。该定理特别包含关于完备诺特局部环结构的Cohen定理（19.8），它们将在第IV章§§6、7中起重要作用。

\noindent IV．研究一个诺特局部环在另一个诺特局部环上的形式光滑性，其意义来自光滑态射的如下“逐点”刻画：若 $X,Y$ 为局部诺特预概形，$f:X\to Y$ 为局部有限型态射，则 $f$ 光滑的充要条件是，对每个 $x\in X$，局部环 $\OO_x$ 在 $\OO_{f(x)}$ 上形式光滑。特别地，若 $Y=\Spec(k)$ 且 $k$ 为完美域，则 $f$ 光滑等价于 $X$ 正则。

\noindent V．最后，我们将在§§20、21、22中看到，形式光滑代数概念自然地出现在Kähler微分理论中，这两种理论彼此阐明。

\egap{19.0.3}
本节及以后各节中，拓扑环与拓扑模均假定为\emph{线性拓扑化}的（$0_{\mathrm I}$，7.1.1）；除非另有明确说明，所考虑的拓扑环均假定交换。回顾：若 $A,B$ 为拓扑环，环同态 $\rho:A\to B$ 定义了 $B$ 的 $A$-代数结构，则称 $B$ 为\emph{拓扑 $A$-代数}，若 $\rho$ 对所考虑的拓扑连续。

为简洁起见，在拓扑环 $A$（或拓扑 $A$-模 $M$）中，将“由理想（或子模）构成的零点基本邻域系”简称为“开理想（或开子模）基本系”。给定拓扑环 $A$ 和一个 $A$-模 $M$，当 $\mathfrak j$ 遍历一个开理想基本系时，子模 $\mathfrak jM$ 构成 $M$ 上某一拓扑的开子模基本系；该拓扑使 $M$ 成为拓扑 $A$-模，称为\emph{由 $A$ 的拓扑诱导的拓扑}。

设拓扑 $A$-模 $M$ 的拓扑比由 $A$ 诱导的拓扑\emph{更粗}。若 $N$ 为 $M$ 的开子模，则离散 $A$-模 $M/N$ 被 $A$ 的某个开理想零化，因为按假设存在开理想 $\mathfrak k$ 使 $\mathfrak kM\subset N$。
若 $M,N$ 的拓扑都由 $A$ 诱导，则每个 $A$-同态 $u:M\to N$ 连续：对 $N$ 中任意零点邻域 $V$，存在开理想 $\mathfrak j$ 使 $\mathfrak jN\subset V$，于是 $u(\mathfrak jM)\subset\mathfrak jN\subset V$。
当 $B$ 为拓扑 $A$-代数时，$B$ 上由 $A$ 诱导的拓扑比其给定拓扑\emph{更细}：对 $B$ 的每个开理想 $\mathfrak k$，按假设存在 $A$ 的开理想 $\mathfrak j$ 使 $\mathfrak jB\subset\mathfrak k$。

\subsection{形式满态射与形式单态射}
\source{印刷第71--78页；上传PDF第70--77页}
\egas{19.1.1}{命题}
设 $A$ 为拓扑环，$M,N$ 为拓扑 $A$-模，$(W_\lambda)$ 为 $N$ 的开子模基本系，$u:M\to N$ 为连续 $A$-同态。以 $\widehat M,\widehat N$ 表示其分离完备化，以 $\widehat u:\widehat M\to\widehat N$ 表示 $u$ 在分离完备化上的连续延拓。

（i）下列条件等价：
\begin{enumerate}[label=\alph*)]
\item $u(M)$ 在 $N$ 中稠密。
\item[$a'$）] $\widehat u(\widehat M)$ 在 $\widehat N$ 中稠密。
\item[$a''$）] 对每个 $\lambda$，复合同态 $M\xrightarrow{u}N\to N/W_\lambda$ 满射。
\end{enumerate}
（ii）下列条件等价：
\begin{enumerate}[label=\alph*),start=2]
\item $N$ 的拓扑经 $u$ 的逆像等于 $M$ 的拓扑。
\item[$b'$）] $\widehat u$ 将拓扑 $\widehat A$-模 $\widehat M$ 同构地映到 $\widehat N$ 的拓扑子 $\widehat A$-模 $\widehat u(\widehat M)$ 上（该子模必为闭子模）。
\item[$b''$）] $u^{-1}(W_\lambda)$ 构成 $M$ 的零点基本邻域系。
\end{enumerate}
这些断言直接来自分离完备化的定义及 $N/W_\lambda$ 的离散性。

\egap{19.1.2}
当（19.1.1）中（i）（或（ii））的等价条件满足时，称 $u$ 为\emph{形式满态射}（或\emph{形式单态射}）。若二者同时满足，则称 $u$ 为\emph{形式双态射}；由（19.1.1），这等价于 $\widehat u$ 是拓扑 $\widehat A$-模 $\widehat M\to\widehat N$ 的同构。\cmt{F6-H02}{这里的“形式”不可删去。例如环 $k[t]$ 配备 $(t)$-预进制拓扑后，到 $k[[t]]$ 的典范映射稠密且诱导原拓扑，但一般并不满射；它是形式双态射而非原模之间的双射。}

\egas{19.1.3}{命题}
设 $A$ 为拓扑环，$M,N$ 为拓扑 $A$-模，$u:M\to N$ 为连续 $A$-同态。假定 $M,N$ 分别有同指标集的开子模基本系 $(V_\lambda)$、$(W_\lambda)$，且 $u(V_\lambda)\subset W_\lambda$。令 $u_\lambda:M/V_\lambda\to N/W_\lambda$ 为商同态。则：

（i）$u$ 为形式满态射，当且仅当每个 $u_\lambda$ 满射。

（ii）$u$ 为形式单态射，当且仅当对每个 $\lambda$，存在 $\mu$ 使 $V_\mu\subset V_\lambda$，且 $\Ker u_\mu$ 包含于典范映射 $M/V_\mu\to M/V_\lambda$ 的核 $V_\lambda/V_\mu$。

（iii）$u$ 为形式双态射，当且仅当对每个 $\lambda$，$u_\lambda$ 满射，且存在 $\mu$ 使 $V_\mu\subset V_\lambda$，并使典范同态 $M/V_\mu\to M/V_\lambda$ 分解为
\[
 M/V_\mu\xrightarrow{u_\mu}N/W_\mu\longrightarrow M/V_\lambda
\]
（其中后一个同态必唯一）。

（i）是直接的。另一方面，$\Ker u_\mu=u^{-1}(W_\mu)/V_\mu$，而典范映射 $M/V_\mu\to M/V_\lambda$ 的核为 $V_\lambda/V_\mu$；所以（ii）的条件等价于 $u^{-1}(W_\mu)\subset V_\lambda$，从而立即得到（ii）。最后，若 $u_\mu$ 满射，则 $N/W_\mu$ 与 $(M/V_\mu)/\Ker u_\mu$ 相同构，因此 $\Ker u_\mu\subset V_\lambda/V_\mu$ 等价于 $M/V_\mu\to M/V_\lambda$ 分解为 $v\circ u_\mu$，其中 $v:N/W_\mu\to M/V_\lambda$ 为同态。

\egas{19.1.4}{推论}
设 $M,N$ 的拓扑均由拓扑环 $A$ 诱导，$u:M\to N$ 为形式满态射，$(\mathfrak j_\lambda)$ 为 $A$ 的开理想基本系。$u$ 是形式双态射的充要条件是，每个商同态 $u_\lambda:M/\mathfrak j_\lambda M\to N/\mathfrak j_\lambda N$ 双射。

事实上，$u(\mathfrak j_\lambda M)\subset\mathfrak j_\lambda N$，故可用（19.1.3，iii）。若有分解
\[
 M/\mathfrak j_\mu M\xrightarrow{u_\mu}N/\mathfrak j_\mu N\xrightarrow{v}M/\mathfrak j_\lambda M,
\]
则 $\mathfrak j_\lambda M/\mathfrak j_\mu M$ 经 $u_\mu$ 的像为 $\mathfrak j_\lambda N/\mathfrak j_\mu N$，所以后者经 $v$ 的像为零，且 $v$ 分解为
\[
 N/\mathfrak j_\mu N\to N/\mathfrak j_\lambda N\xrightarrow{w}M/\mathfrak j_\lambda M.
\]
于是 $M/\mathfrak j_\lambda M$ 的恒等自同构分解为
\[
 M/\mathfrak j_\lambda M\xrightarrow{u_\lambda}N/\mathfrak j_\lambda N\xrightarrow{w}M/\mathfrak j_\lambda M.
\]
这表明 $u_\lambda$ 单射；它已经是满射，推论得证。

\egas{19.1.5}{命题}
设 $A$ 为拓扑环，$M,N$ 为拓扑 $A$-模，$u:M\to N$ 为连续 $A$-同态。下列条件等价：

$a)$ 对每个离散拓扑 $A$-模 $E$ 和每个连续 $A$-同态 $v:M\to E$，存在连续 $A$-同态 $w:N\to E$ 使 $v=w\circ u$。

$b)$ 对 $M$ 的每个开子模 $V'$，存在 $N$ 的开子模 $W''$、$M$ 的开子模 $V''\subset V'\cap u^{-1}(W'')$，以及同态 $h:N/W''\to M/V'$，使典范同态 $M/V''\to M/V'$ 分解为
\[
 M/V''\xrightarrow{u''}N/W''\xrightarrow{h}M/V',
\]
其中 $u''$ 由 $u$ 取商得到。

要由 $a)$ 推出 $b)$，只须将 $a)$ 用于 $E=M/V'$ 和典范同态 $v:M\to M/V'$。这时 $W''=w^{-1}(0)$ 是 $N$ 的开子模，$V''=u^{-1}(W'')$ 是包含于 $V'$ 的开子模；这些子模和由 $w$ 取商所得的 $h:N/W''\to M/V'$ 满足要求。
反之，将 $E$ 换成 $v(M)$，可在证明 $b)\Rightarrow a)$ 时假定 $v$ 满射。于是 $V'=\Ker v$ 是开子模，$E$ 同构于离散 $A$-模 $M/V'$。应用 $b)$，取复合 $N\to N/W''\xrightarrow{h}M/V'$ 即得所求连续同态 $w$，因为下图交换：
\[
\begin{tikzcd}
 M\arrow[r,"u"]\arrow[d] & N\arrow[d]\\
 M/V''\arrow[r,"u''"] & N/W''.
\end{tikzcd}
\]
当这些条件满足时，称 $u$ \emph{形式可左逆}。由于 $b)$ 蕴含 $u^{-1}(W'')\subset V'$，$u$ 此时为形式单态射。以下推论也说明该术语的合理性。

\egas{19.1.6}{推论}
若存在连续 $\widehat A$-同态 $s:\widehat N\to\widehat M$，满足 $s\circ\widehat u=1_{\widehat M}$，则 $u$ 形式可左逆。

注意，此时 $\widehat u$ 将 $\widehat M$ 拓扑同构到 $\widehat N$ 的一个拓扑直和因子上。沿用（19.1.5，$a$）的记号，因 $E$ 离散，$v:M\to E$ 可延拓为连续 $\widehat A$-同态 $\widehat v:\widehat M\to E$。设 $j:M\to\widehat M$、$j':N\to\widehat N$ 为典范同态；令 $w=\widehat v\circ s\circ j'$，则
\[
 w\circ u=\widehat v\circ s\circ\widehat u\circ j=\widehat v\circ j=v.
\]

\egas{19.1.7}{推论}
设 $M,N$ 的拓扑均由 $A$ 诱导，$(\mathfrak j_\lambda)$ 为 $A$ 的开理想基本系。$u$ 形式可左逆，当且仅当每个商同态
$u_\lambda:M/\mathfrak j_\lambda M\to N/\mathfrak j_\lambda N$ 可左逆，即它将 $M/\mathfrak j_\lambda M$ 同构到 $N/\mathfrak j_\lambda N$ 的一个直和因子上。

充分性：在（19.1.5，$b$）中取 $V'=\mathfrak j_\lambda M$、$W''=\mathfrak j_\lambda N$、$V''=V'$，以及满足 $h\circ u_\lambda=1$ 的 $h$。反之，若 $u$ 形式可左逆，则对每个 $\lambda$，由（19.1.5，$b$）存在 $\mathfrak j_\mu\subset\mathfrak j_\lambda$ 及 $h:N/\mathfrak j_\mu N\to M/\mathfrak j_\lambda M$，使典范同态分解为
\[
 M/\mathfrak j_\mu M\xrightarrow{u_\mu}N/\mathfrak j_\mu N\xrightarrow{h}M/\mathfrak j_\lambda M.
\]
由于
\[
 h\bigl(\mathfrak j_\lambda(N/\mathfrak j_\mu N)\bigr)
 =\mathfrak j_\lambda h(N/\mathfrak j_\mu N)
 \subset\mathfrak j_\lambda(M/\mathfrak j_\lambda M)=0,
\]
$h$ 典范地分解为 $N/\mathfrak j_\mu N\to N/\mathfrak j_\lambda N\xrightarrow{s}M/\mathfrak j_\lambda M$，且 $s$ 是 $u_\lambda$ 的左逆。

\egas{19.1.8}{命题}
设拓扑环 $A$ 有可数递降开理想基本系 $(\mathfrak j_n)$；$M,N$ 为拓扑由 $A$ 诱导的拓扑 $A$-模，$u:M\to N$ 为 $A$-同态。令 $M_n=M/\mathfrak j_nM$、$N_n=N/\mathfrak j_nN$，$u_n:M_n\to N_n$ 为商同态。假定每个 $P_n=\Coker u_n$ 为投射 $(A/\mathfrak j_n)$-模，且 $M$ 分离完备。则 $u$ 形式可左逆的充要条件是 $u$ 可左逆；此时 $u$ 将 $M$ 拓扑同构到 $N$ 的一个拓扑直和因子上。

由（19.1.7），有交换图
\[
\begin{tikzcd}[column sep=1.8em]
 0\arrow[r]&M_n\arrow[r,"u_n"]&N_n\arrow[r,"p_n"]&P_n\arrow[r]&0\\
 0\arrow[r]&M_{n+1}\arrow[u,"f"]\arrow[r,"u_{n+1}"]&N_{n+1}\arrow[u,"g"]\arrow[r,"p_{n+1}"]&P_{n+1}\arrow[u,"h"]\arrow[r]&0,
\end{tikzcd}
\]
各行是分裂正合列，亦即对每个 $n$ 存在 $s'_n:P_n\to N_n$ 使 $p_n\circ s'_n=1_{P_n}$。我们将对 $n$ 递归构造 $s_n:P_n\to N_n$，满足 $p_n\circ s_n=1_{P_n}$，且
\[
\begin{tikzcd}
 P_n\arrow[r,"s_n"]&N_n\\
 P_{n+1}\arrow[u,"h"]\arrow[r,"s_{n+1}"]&N_{n+1}\arrow[u,"g"]
\end{tikzcd}
\]
交换。事实上，$g\circ s'_{n+1}-s_n\circ h$ 是从 $P_{n+1}$ 到
\[
 u_n(M_n)=u_{n+1}(M_{n+1})/\mathfrak j_nu_{n+1}(M_{n+1})
\]
的同态；由于 $P_{n+1}$ 为投射 $(A/\mathfrak j_{n+1})$-模，该同态分解为
\[
 P_{n+1}\xrightarrow{t_{n+1}}u_{n+1}(M_{n+1})
 \longrightarrow u_{n+1}(M_{n+1})/\mathfrak j_nu_{n+1}(M_{n+1}).
\]
因此 $s_{n+1}=s'_{n+1}-t_{n+1}$ 满足要求。由 $N_n=u_n(M_n)\oplus s_n(P_n)$，得到 $u_n$ 的左逆 $w_n:N_n\to M_n$，而且下图交换：
\[
\begin{tikzcd}
 N_n\arrow[r,"w_n"]&M_n\\
 N_{n+1}\arrow[u,"g"]\arrow[r,"w_{n+1}"]&M_{n+1}\arrow[u,"f"].
\end{tikzcd}
\]
射影系 $(w_n)$ 有射影极限 $\widehat w:\widehat N\to\widehat M=M$。再与 $N\to\widehat N$ 复合，得 $w:N\to M$，使 $w\circ u$ 在每个 $M_n$ 上诱导恒等映射。$M$ 分离完备，故 $w\circ u=1_M$。

\egas{19.1.9}{命题}
设 $A$ 为预可容许拓扑环（$0_{\mathrm I}$，7.1.2），$\mathfrak a$ 为定义理想，$(\mathfrak j_\lambda)$ 为开理想基本系。设 $M,N$ 的拓扑由 $A$ 诱导，且每个 $N/\mathfrak j_\lambda N$ 为投射 $(A/\mathfrak j_\lambda)$-模（参见19.2.3）。若 $u:M\to N$ 为 $A$-同态，则下列条件等价：

$a)$ $u$ 形式可左逆。

$b)$ 商同态 $u_0:M/\mathfrak aM\to N/\mathfrak aN$ 可左逆。

由（19.1.7），$a)$ 等价于每个 $u_\lambda$ 可左逆。因 $\mathfrak a$ 开，故包含某个 $\mathfrak j_\lambda$，从而 $u_0$ 可左逆。反过来，为证 $b)\Rightarrow a)$，注意对每个 $\lambda$，$\mathfrak a/\mathfrak j_\lambda$ 按假设为 $A/\mathfrak j_\lambda$ 的幂零理想。\cmt{F6-T02}{此处须将开理想基本系限于包含在 $\mathfrak a$ 中的共尾子系；否则 $\mathfrak a/\mathfrak j_\lambda$ 未必有定义。也可统一写为 $(\mathfrak a+\mathfrak j_\lambda)/\mathfrak j_\lambda$。这是证明中可直接实行的基本系收缩，不是增加命题的实质假设。}所需断言由下命题推出。

\egas{19.1.10}{命题}
设 $A$ 为环，$M,N$ 为 $A$-模，$N$ 投射，$u:M\to N$ 为 $A$-同态。设理想 $\mathfrak j$ 满足下列条件之一：

（i）$\mathfrak j$ 幂零；

（ii）$\mathfrak j$ 包含于 $A$ 的根中，且 $M$ 有限型。

则 $u$ 可左逆，当且仅当商同态 $u_0:M/\mathfrak jM\to N/\mathfrak jN$ 可左逆。

必要性是直接的。为证充分性，设 $v_0$ 为 $u_0$ 的左逆。因 $N$ 投射，复合 $N\to N/\mathfrak jN\xrightarrow{v_0}M/\mathfrak jM$ 分解为 $N\xrightarrow{v}M\to M/\mathfrak jM$。于是 $w=v\circ u$ 在 $M/\mathfrak jM$ 上诱导恒等自同构 $w_0$。只须证明 $w$ 双射，因为此时 $w^{-1}\circ v$ 即为 $u$ 的左逆。分两种情况。

（i）对每个 $n$ 有交换图
\[
\begin{tikzcd}[column sep=large]
 (\mathfrak j^n/\mathfrak j^{n+1})\otimes_{A/\mathfrak j}(M/\mathfrak jM)
 \arrow[r]\arrow[d,"1\otimes\gr^0(w)"']&\mathfrak j^nM/\mathfrak j^{n+1}M\arrow[d,"\gr^n(w)"]\\
 (\mathfrak j^n/\mathfrak j^{n+1})\otimes_{A/\mathfrak j}(M/\mathfrak jM)
 \arrow[r]&\mathfrak j^nM/\mathfrak j^{n+1}M.
\end{tikzcd}
\]
横箭头满射，而 $\gr^0(w)=w_0$ 为恒等映射，故 $\gr^n(w)$ 也为恒等映射，特别是双射。因 $\mathfrak j$ 幂零，$M$ 的 $\mathfrak j$-预进制滤过有限；所以 $w$ 双射（Bourbaki，\textit{交换代数}，第三章，§2，第8目，定理1的推论3）。

（ii）只须对每个极大理想 $\mathfrak m$ 证明 $w_{\mathfrak m}$ 在 $M_{\mathfrak m}$ 上双射（同书，第二章，§3，第3目，定理1）。由于 $\mathfrak jA_{\mathfrak m}\subset\mathfrak mA_{\mathfrak m}$，且 $A_{\mathfrak m}/\mathfrak jA_{\mathfrak m}=(A/\mathfrak j)_{\mathfrak m}$，可归结为局部环。此外可假定 $\mathfrak j$ 为极大理想 $\mathfrak j_0$：若 $u_0$ 可左逆，则它与 $1_{A/\mathfrak j_0}$ 张量所得的 $u_{00}:M/\mathfrak j_0M\to N/\mathfrak j_0N$ 也可左逆，因为 $M/\mathfrak j_0M=(M/\mathfrak jM)\otimes_{A/\mathfrak j}(A/\mathfrak j_0)$。

故设 $\mathfrak j$ 极大，$A/\mathfrak j$ 为域。只须证明 $M$ 在给定条件下是自由 $A$-模：这时 $\det w_0$ 是 $\det w$ 在 $A/\mathfrak j$ 中的像，故 $\det w\notin\mathfrak j$，从而可逆。因 $M/\mathfrak jM$ 为有限维向量空间，可取有限型自由 $A$-模 $L$ 及 $f:L\to M$，使商同态 $f_0:L/\mathfrak jL\to M/\mathfrak jM$ 双射。\cmt{F6-T03}{初版此处仅印“有限型 $A$-模 $L$”，紧接着却使用“这里 $L$ 是自由的”。应在选择 $L$ 时明确取有限型自由模：提升 $M/\mathfrak jM$ 的一组基即可。此处在正文补出“自由”，不改变原证明的方法。}由Nakayama引理，$f$ 满射（同书，第二章，§3，第2目，命题4的推论1）。令 $g=u\circ f$，则 $g_0:L/\mathfrak jL\to N/\mathfrak jN$ 可左逆。由于 $L$ 自由，开头的论证证明 $g$ 本身可左逆；这迫使 $f$ 单射，证明结束。

顺便记录（19.1.10）的如下推论。

\egas{19.1.11}{推论}
设 $A$ 为局部环，$k$ 为其剩余域，$M$ 为有限型 $A$-模，$N$ 为投射 $A$-模，$u:M\to N$ 为同态。$u$ 可左逆，当且仅当 $M$ 有一个生成系 $(x_i)_{1\leq i\leq m}$，使 $x_i\otimes1$ 经 $u\otimes1:M\otimes_Ak\to N\otimes_Ak$ 的像线性无关；此时 $x_i$ 是 $M$ 的一组基。

若 $u$ 可左逆，则 $M$ 为有限型投射模，因而自由，必要性成立。反之，条件满足时，$x_i\otimes1$ 是 $M\otimes_Ak$ 的一组基，且 $u\otimes1$ 可左逆。对极大理想 $\mathfrak j$ 应用（19.1.10）即可。

\egas{19.1.12}{推论}
设 $A$ 为环，$M$ 为有限型 $A$-模，$N$ 为投射 $A$-模，$u:M\to N$ 为同态。对每个 $\mathfrak p\in\Spec(A)$，下列条件等价：

$a)$ $u_{\mathfrak p}:M_{\mathfrak p}\to N_{\mathfrak p}$ 可左逆。

$b)$ $u\otimes1:M\otimes_Ak(\mathfrak p)\to N\otimes_Ak(\mathfrak p)$ 单射，其中 $k(\mathfrak p)$ 为 $A$ 在 $\mathfrak p$ 处的剩余域。

$c)$ 存在有限个 $x_i\in M$（$1\leq i\leq m$），其在 $M_{\mathfrak p}$ 中的像生成该模，并存在 $N$ 上的 $m$ 个线性形式 $v_i$，使
$\det(\langle v_i,u(x_j)\rangle)\notin\mathfrak p$。

$d)$ 存在 $f\in A\setminus\mathfrak p$，使 $u_f:M_f\to N_f$ 可左逆。

此外，满足这些条件的素理想构成 $\Spec(A)$ 的开集。

最后一项断言直接由 $d)$ 得出。因为 $N$ 投射，它是某个自由模 $A^{(I)}$ 的直和因子；又因 $M$ 有限型，$u(M)$ 包含于一个有限秩自由坐标子模 $A^n$。将 $N$ 换成 $N\oplus P$（$u$ 的像仍在 $N$ 内），$a)$、$b)$、$c)$、$d)$ 各自均保持等价，故可限于 $N$ 有限型自由的情形。
$d)\Rightarrow a)$ 是直接的；$a)$ 与 $b)$ 由（19.1.11）等价。$a)$ 还蕴含 $M_{\mathfrak p}$ 自由，因此可取 $x_i$ 的像为其一组基；注意 $N$ 有限型自由，$N_{\mathfrak p}$ 上的任意线性形式可写成 $v/f$，其中 $f\notin\mathfrak p$、$v$ 为 $N$ 上的线性形式，便得 $c)$。$c)\Rightarrow b)$ 也是直接的。

只剩 $a)\Rightarrow d)$。因 $N$ 有限表示，有典范同构
\[
 (\Hom_A(N,M))_{\mathfrak p}\simeq
 \Hom_{A_{\mathfrak p}}(N_{\mathfrak p},M_{\mathfrak p})
\]
（Bourbaki，\textit{交换代数}，第二章，§2，第7目，命题19）。若 $w_{\mathfrak p}$ 为 $u_{\mathfrak p}$ 的左逆，则存在 $w:N\to M$ 和 $f\notin\mathfrak p$，使 $w_{\mathfrak p}$ 是 $w$ 在局部化后除以 $f$ 所得的同态。\cmt{F6-T04}{初版此处将 $w$ 与 $w_{\mathfrak p}$ 的等式方向排混；本句按随后的关系 $(w\circ u)_{\mathfrak p}=f\,1_{M_{\mathfrak p}}$ 明示记号，以免把局部同态与全局同态写成无类型的等式。}于是 $(w\circ u)_{\mathfrak p}=f\,1_{M_{\mathfrak p}}$。由 $M$ 有限型，存在 $g\notin\mathfrak p$，使 $g(w\circ u-f1_M)$ 在所有生成元上为零，故为零。令 $h=fg$，则在 $A_h$ 上 $g$、$f$ 均可逆，并有 $w_h\circ u_h=(f/1)1_{M_h}$。所以 $(f/1)^{-1}w_h$ 是 $u_h$ 的左逆。

\egas{19.1.13}{命题}
在（19.1.9）的假设外，再假定开理想基本系为递降序列 $(\mathfrak j_n)$，且 $M$ 分离完备。那么（19.1.9）的条件 $a)$、$b)$ 还等价于：

$c)$ $u$ 可左逆。

已知 $c)\Rightarrow a)$（19.1.6）。反过来，若 $a)$ 成立，沿用（19.1.8）的记号，$M_n$ 是投射 $(A/\mathfrak j_n)$-模 $N_n$ 的直和因子，故 $P_n$ 也同构于 $N_n$ 的直和因子，因而投射；应用（19.1.8）即可。

最后记录以下命题。

\egas{19.1.14}{命题}
设 $A$ 为环，$M$ 为有限型 $A$-模，$N$ 为投射 $A$-模，$u:M\to N$ 为同态。

（i）$u$ 可左逆，当且仅当对 $A$ 的每个极大理想 $\mathfrak m$，$u_{\mathfrak m}:M_{\mathfrak m}\to N_{\mathfrak m}$ 可左逆。

（ii）设 $A'$ 为忠实平坦的 $A$-代数。$u$ 可左逆，当且仅当 $u\otimes1:M\otimes_AA'\to N\otimes_AA'$ 可左逆。

如（19.1.12），可设 $N$ 有限型自由。于是说 $u$ 可左逆，等价于 $u$ 单射且商模 $P=N/u(M)$ 投射，因为这时 $u(M)$ 为 $N$ 的直和因子。注意 $M$ 有限型，所以 $P$ 有限表示。

（i）必要性是直接的。反之，条件成立时，$u$ 单射（Bourbaki，\textit{交换代数}，第二章，§3，第3目，定理1）；而每个 $P_{\mathfrak m}=N_{\mathfrak m}/u_{\mathfrak m}(M_{\mathfrak m})$ 投射，故 $P$ 投射（同书，§5，第2目，定理1）。

（ii）必要性仍是直接的。反之，$u$ 单射（$0_{\mathrm I}$，6.4.1），且 $P\otimes_AA'=\Coker(u\otimes1)$ 投射，因而平坦。故 $P$ 为平坦 $A$-模（$0_{\mathrm I}$，6.6.3）；它又有限表示，所以投射（Bourbaki，\textit{交换代数}，第二章，§5，第2目，定理1的推论2）。

\egas{19.1.15}{注记}
形式可左逆同态的“对偶”概念是\emph{形式可右逆}同态。按定义，连续 $A$-同态 $u:M\to N$ 满足如下条件：对 $N$ 的每个开子模 $W'$，存在 $M$ 的开子模 $V'\subset u^{-1}(W')$、$N$ 的开子模 $W''\subset W'$，以及 $h:N/W''\to M/V'$，使典范同态 $N/W''\to N/W'$ 分解为
\[
 N/W''\xrightarrow{h}M/V'\xrightarrow{u'}N/W',
\]
其中 $u'$ 由 $u$ 取商得到。这蕴含 $u$ 为形式满态射。若存在连续 $\widehat A$-同态 $r:\widehat N\to\widehat M$，使 $\widehat u\circ r=1_{\widehat N}$，则可直接验证 $u$ 形式可右逆。
若（19.1.7）的条件成立，$u$ 形式可右逆当且仅当每个 $u_\lambda$ 可右逆，即 $\Ker u_\lambda$ 为 $M/\mathfrak j_\lambda M$ 的直和因子，且 $u_\lambda$ 将该核的某个补模同构到 $N/\mathfrak j_\lambda N$。
与（19.1.8）、（19.1.9）对应的命题留给读者表述并证明：在（19.1.8）的对应命题中，须假定 $M$ 分离完备且 $M_n$ 为投射 $(A/\mathfrak j_n)$-模；在（19.1.9）的对应命题中，应删去关于 $N/\mathfrak j_\lambda N$ 的假设，改为假定每个 $M/\mathfrak j_\lambda M$ 为投射 $(A/\mathfrak j_\lambda)$-模。
